% !TeX root = ./rpz.tex
\setcounter{chapter}{6}
\appendixchapter{В}{Техническое задание на разработку}{cha:appendix_tz}

\begin{center}(обязательное)\end{center}

\section*{В.1 Общие сведения}

Полное наименование системы: веб-приложение для автоматизированного планирования маршрутов с учетом временных окон.

Условное обозначение: <<Планировщик маршрутов>>.

Наименование организации-разработчика: Московский политехнический университет.

Основание для разработки: задание на выпускную квалификационную работу.

Плановые сроки начала и окончания работы: февраль 2026~\textendash{} май 2026~года.

Источники финансирования: разработка ведется в рамках выпускной квалификационной работы и не предусматривает прямого финансирования.

Порядок оформления и предъявления заказчику результатов работы: результаты передаются заказчику в виде комплекта документации (пояснительная записка), исходного кода (репозиторий Git) и собранного дистрибутива (Docker-образ серверной части, статическая сборка клиентской части).

\section*{В.2 Назначение и цели создания системы}

Назначение системы. Система предназначена для автоматизированного построения оптимальной последовательности выполнения задач с учетом географического положения точек, временных окон обслуживания и длительности выполнения каждой задачи.

Цели создания системы:

\begin{compactlist}
  \item сокращение времени, затрачиваемого пользователем на ручное планирование маршрута с временными окнами;
  \item уменьшение числа ошибок планирования (опоздания, нарушения временных окон) за счет автоматизированной проверки допустимости;
  \item обеспечение доступности инструментов оптимизации маршрута для частных пользователей и малого бизнеса без необходимости приобретения коммерческой подписки.
\end{compactlist}

\section*{В.3 Характеристика объекта автоматизации}

Объект автоматизации: процесс планирования последовательности выполнения задач пользователя в городской среде.

Существующий процесс. В отсутствие специализированного программного обеспечения пользователь выполняет планирование вручную: оценивает время перемещения между адресами по сторонним картографическим сервисам, учитывает временные окна и длительность выполнения, формирует порядок посещения. При увеличении числа задач трудоемкость процесса возрастает нелинейно, а вероятность нарушения временных ограничений увеличивается.

Целевая аудитория: частные пользователи, специалисты на выезде, курьеры малого бизнеса, индивидуальные предприниматели, диспетчеры, формирующие маршруты вручную.

\section*{В.4 Требования к системе}

\subsection*{В.4.1 Требования к системе в целом}

Структура системы. Система реализуется в виде клиент-серверного веб-приложения и состоит из двух независимых компонентов:

\begin{compactlist}
  \item серверной части, реализующей хранение данных, бизнес-логику и алгоритмы оптимизации;
  \item клиентской части, обеспечивающей взаимодействие с пользователем через веб-браузер.
\end{compactlist}

Численность пользователей: до 100 одновременных активных пользователей.

Режим функционирования: непрерывный (24/7) с допустимыми технологическими перерывами на обслуживание и обновление.

Показатели назначения. Время отклика интерфейса при типовом сценарии не превышает 2~секунд при нагрузке до 5 запросов в секунду на сервер.

Требования к надежности. Система должна сохранять работоспособность при кратковременной недоступности внешнего картографического сервиса за счет механизма кэширования.

Требования к эргономике. Пользовательский интерфейс должен быть адаптивным и сохранять работоспособность на устройствах с диагональю экрана от 4 дюймов (смартфон) до 32 дюймов (настольный монитор).

Требования к безопасности информации. Пароли пользователей хранятся в виде хэш-значений (алгоритм \texttt{bcrypt}); токены обновления сессии~\textendash{} в виде хэш-значений алгоритма SHA-256. Передача данных между клиентом и сервером в продуктивной среде выполняется по протоколу HTTPS.

\subsection*{В.4.2 Требования к функциям}

Система должна реализовать следующие функции:

\begin{compactlist}
  \item регистрация и аутентификация пользователей по электронной почте и паролю;
  \item создание, редактирование, удаление и клонирование задач с указанием адреса, длительности выполнения и временного окна;
  \item импорт перечня задач из файлов формата CSV и XLSX;
  \item экспорт перечня задач в форматы CSV и XLSX;
  \item геокодирование адресов через внешний сервис картографии;
  \item построение матрицы расстояний между адресами задач;
  \item автоматизированное вычисление оптимальной последовательности выполнения задач с учетом временных окон;
  \item поддержка дополнительных ограничений: фиксация начальной и конечной точки маршрута, попарные ограничения порядка выполнения задач;
  \item расчет времени прибытия, ожидания и завершения обслуживания для каждой остановки;
  \item визуализация маршрута на карте с указанием порядка посещения;
  \item обнаружение и подсветка конфликтов временных окон;
  \item выбор способа перемещения (автомобиль, пешеход, общественный транспорт).
\end{compactlist}

\subsection*{В.4.3 Требования к видам обеспечения}

Требования к информационному обеспечению. В качестве СУБД используется PostgreSQL версии~16. Структура базы данных управляется механизмом миграций.

Требования к программному обеспечению серверной части: операционная система с поддержкой контейнеризации (Linux, macOS, Windows); Docker Engine версии 24.0 и выше; Docker Compose версии 2.0 и выше.

Требования к программному обеспечению клиентской части: современный веб-браузер с поддержкой JavaScript ES2020 (Google Chrome 90+, Mozilla Firefox 88+, Safari 14+, Microsoft Edge 90+, Yandex Browser 25+).

Требования к техническому обеспечению серверной части: процессор с тактовой частотой не ниже 1~ГГц; оперативная память не менее 512~МБ; дисковое пространство не менее 1~ГБ; постоянное соединение с сетью Интернет.

Требования к лингвистическому обеспечению: основной язык пользовательского интерфейса~\textendash{} русский.

\section*{В.5 Состав и содержание работ по созданию системы}

Разработка системы выполняется в следующей последовательности:

\begin{compactlist}
  \item анализ предметной области, постановка задачи, формирование требований;
  \item проектирование архитектуры программной системы, модели данных;
  \item исследование и выбор алгоритма оптимизации маршрута;
  \item реализация серверной части (REST API, бизнес-логика, репозитории);
  \item реализация клиентской части (пользовательский интерфейс, интеграция с картографией);
  \item модульное тестирование, нагрузочное тестирование;
  \item подготовка эксплуатационной документации.
\end{compactlist}

\section*{В.6 Порядок контроля и приемки}

Приемка результатов работы осуществляется в форме защиты выпускной квалификационной работы перед государственной экзаменационной комиссией. Приемочные испытания включают демонстрацию работы системы на типовом сценарии, проверку выполнения функциональных и нефункциональных требований, а также проверку отчета о нагрузочном тестировании.

\section*{В.7 Требования к составу и содержанию работ по подготовке объекта автоматизации}

Для ввода системы в действие на стороне эксплуатирующей организации необходимо:

\begin{compactlist}
  \item обеспечить наличие сервера с установленным Docker Engine версии 24.0 и выше и Docker Compose версии 2.0 и выше;
  \item получить ключ доступа к Яндекс JavaScript API на портале разработчиков;
  \item установить значение секретного ключа для подписи JWT-токенов длиной не менее 32~символов;
  \item сконфигурировать список допустимых источников CORS;
  \item обеспечить резервное копирование данных PostgreSQL средствами штатной утилиты \texttt{pg\_dump}.
\end{compactlist}

\section*{В.8 Требования к документированию}

Состав поставляемой документации:

\begin{compactlist}
  \item пояснительная записка к выпускной квалификационной работе;
  \item исходный код проекта с комментариями в формате \texttt{godoc} (серверная часть) и \texttt{TSDoc} (клиентская часть);
  \item файл \texttt{README.md} с инструкцией по развертыванию;
  \item руководство пользователя и администратора (раздел~2 пояснительной записки).
\end{compactlist}

\section*{В.9 Источники разработки}

При разработке системы использованы следующие источники:

\begin{compactlist}
  \item ГОСТ 34.602-89. Информационная технология. Комплекс стандартов на автоматизированные системы. Техническое задание на создание автоматизированной системы;
  \item ГОСТ 7.32-2017. Система стандартов по информации, библиотечному и издательскому делу. Отчет о научно-исследовательской работе. Структура и правила оформления;
  \item ГОСТ 19.701-90 (ИСО 5807-85). Единая система программной документации. Схемы алгоритмов, программ, данных и систем;
  \item публикации по теории графов и задачам комбинаторной оптимизации, приведенные в списке использованных источников пояснительной записки;
  \item официальная документация выбранных программных платформ (Go, PostgreSQL, React, Vite, Яндекс JS API).
\end{compactlist}
